% !TeX program = pdflatex
\documentclass[aspectratio=169,xcolor={dvipsnames,table}]{beamer}

% Preámbulo modular para presentaciones Beamer (Modelación Estadística)
% Diseño y paleta institucional del Tecnológico de Monterrey

\documentclass[aspectratio=169,xcolor={dvipsnames,table}]{beamer}

\usepackage[utf8]{inputenc}
\usepackage[spanish,mexico]{babel}
\usepackage{amsmath,amssymb,amsfonts,mathtools}
\usepackage{graphicx}
\usepackage{listings}
\usepackage{booktabs}
\usepackage{tikz}
\usetikzlibrary{matrix,arrows,decorations.pathmorphing,shapes.geometric,positioning}

% Paleta Institucional Tecnológico de Monterrey
\definecolor{TecAzul}{HTML}{0039A6}       % Azul TEC: color primario institucional
\definecolor{TecAzulOscuro}{HTML}{1A2E51} % Azul marino oscuro
\definecolor{TecRojo}{HTML}{EC2661}       % Rojo de acento (paleta miTec)
\definecolor{TecGris}{HTML}{646464}       % Gris neutro para comentarios y subtítulos
\definecolor{TecFondoBloque}{HTML}{F4F6F9}% Fondo suave para bloques

% Configuración de Tema Beamer
\usetheme{Boadilla}
\usecolortheme[named=TecAzulOscuro]{structure}
\setbeamercolor{alerted text}{fg=TecRojo}
\setbeamercolor{palette primary}{bg=TecAzulOscuro,fg=white}
\setbeamercolor{palette secondary}{bg=TecAzul,fg=white}
\setbeamercolor{palette tertiary}{bg=TecAzulOscuro!80!black,fg=white}
\setbeamercolor{title}{fg=TecAzulOscuro,bg=TecFondoBloque}
\setbeamercolor{frametitle}{fg=TecAzulOscuro,bg=TecFondoBloque}

% Bloques personalizados elegantes
\setbeamercolor{block title}{bg=TecAzulOscuro,fg=white}
\setbeamercolor{block body}{bg=TecFondoBloque,fg=black}
\setbeamercolor{block title alerted}{bg=TecRojo,fg=white}
\setbeamercolor{block body alerted}{bg=TecRojo!8,fg=black}
\setbeamercolor{block title example}{bg=TecAzul,fg=white}
\setbeamercolor{block body example}{bg=TecAzul!8,fg=black}

% Redondear bordes y añadir sombra sutil
\setbeamertemplate{blocks}[rounded][shadow=true]
\setbeamertemplate{navigation symbols}{}

% Pie de página limpio con número de diapositiva
\setbeamertemplate{footline}{
  \leavevmode%
  \hbox{%
  \begin{beamercolorbox}[wd=.7\paperwidth,ht=2.25ex,dp=1ex,left]{author in head/foot}%
    \hspace*{2ex}\insertshortauthor~~(\insertshortinstitute) \hfill \insertshorttitle
  \end{beamercolorbox}%
  \begin{beamercolorbox}[wd=.3\paperwidth,ht=2.25ex,dp=1ex,right]{date in head/foot}%
    \insertshortdate{}\hspace*{2em}
    \insertframenumber{} / \inserttotalframenumber\hspace*{2ex}
  \end{beamercolorbox}}%
  \vskip0pt%
}

% Configuración de Listings de Python para visualización pedagógica de código
\lstset{
    language=Python,
    basicstyle=\ttfamily\footnotesize,
    keywordstyle=\color{TecAzulOscuro}\bfseries,
    stringstyle=\color{TecRojo},
    commentstyle=\color{TecGris}\itshape,
    showstringspaces=false,
    breaklines=true,
    frame=lines,
    rulecolor=\color{TecAzulOscuro!30},
    backgroundcolor=\color{TecFondoBloque},
    aboveskip=0.5em,
    belowskip=0.5em,
    literate={á}{{\'a}}1 {é}{{\'e}}1 {í}{{\'i}}1 {ó}{{\'o}}1 {ú}{{\'u}}1
             {Á}{{\'A}}1 {É}{{\'E}}1 {Í}{{\'I}}1 {Ó}{{\'O}}1 {Ú}{{\'U}}1
             {ñ}{{\~n}}1 {Ñ}{{\~N}}1 {¿}{{?`}}1 {¡}{{!`}}1
}

% Metadatos institucionales por defecto
\title[Modelación Estadística]{Modelación Estadística para Ciencia de Datos e Ingeniería}
\author[J. Castillo Colmenares]{Juliho David Castillo Colmenares}
\institute[Tec de Monterrey]{Tecnológico de Monterrey \\ Escuela de Ingeniería y Ciencias}
\date{\today}

% Comandos y macros estadísticas compartidas para las presentaciones Beamer (Metropolis)

% Conjuntos numéricos básicos
\newcommand{\R}{\mathbb{R}}
\newcommand{\N}{\mathbb{N}}
\newcommand{\Q}{\mathbb{Q}}
\newcommand{\Z}{\mathbb{Z}}
\newcommand{\set}[1]{\left\{ #1 \right\}}
\newcommand{\abs}[1]{\left|#1\right|}
\newcommand{\norm}[1]{\left\| #1 \right\|}

% Comandos estadísticos y probabilísticos del curso
\newcommand{\Var}{\operatorname{Var}}
\newcommand{\cov}{\operatorname{Cov}}
\newcommand{\corr}{\rho}
\newcommand{\comb}[2]{\begin{pmatrix} #1 \\ #2 \end{pmatrix}}
\newcommand{\s}{\sigma}
\newcommand{\particion}{\mathfrak{P}}
\newcommand{\card}[1]{\#\left( #1 \right)}
\newcommand{\seq}[3]{\set{#1}_{#2}^{#3}}
\newcommand{\E}{\mathbb{E}}
\newcommand{\Prob}{\mathbb{P}}

% Entornos pedagógicos y cajas en español académico natural
\newenvironment{discusion}{%
  \begin{alertblock}{Pregunta de Discusión}%
}{%
  \end{alertblock}%
}

\newenvironment{reflexion}{%
  \begin{alertblock}{Reflexión para el Aula}%
}{%
  \end{alertblock}%
}

\newenvironment{paradoja}{%
  \begin{alertblock}{Paradoja de Exploración}%
}{%
  \end{alertblock}%
}

\newenvironment{motivacion}{%
  \begin{exampleblock}{Motivación: Ciencia de Datos en la Práctica}%
}{%
  \end{exampleblock}%
}

\newenvironment{retoclase}{%
  \begin{block}{Reto Rápido en Equipo}%
}{%
  \end{block}%
}


\title{Funciones de Probabilidad Discretas}
\subtitle{Sección 03.01 --- Conceptos Básicos, Soporte y PMF}
\author[J. Castillo Colmenares]{Juliho Castillo Colmenares}
\institute[Tec de Monterrey]{Tecnológico de Monterrey}
\date{}

\begin{document}

% SLIDE 01: Portada
\begin{frame}[plain]
	\titlepage
\end{frame}

% SLIDE 02: Hoja de Ruta
\begin{frame}{Hoja de Ruta --- Capítulo 03: Variables Aleatorias Discretas}
	\begin{block}{Estructura Curricular de la Unidad 2}
		\small
		\begin{enumerate}
			\itemsep=0.04cm
			\item \textbf{\textcolor{TecRojo}{Sección 03.01: Funciones de Probabilidad Discretas (PMF y Soporte)}} $\leftarrow$ \emph{Foco actual}
			\item \textcolor{gray}{Sección 03.02: Función de Distribución Acumulada (CDF Discreta)}
			\item \textcolor{gray}{Sección 03.03: Esperanza Matemática, Varianza y Momentos}
			\item \textcolor{gray}{Sección 03.04: Distribuciones Especiales I (Binomial, Geométrica e Hipergeométrica)}
			\item \textcolor{gray}{Sección 03.05: Distribuciones Especiales II (Poisson, Binomial Negativa y Multinomial)}
		\end{enumerate}
	\end{block}
	\vspace{-0.15cm}
	\begin{alertblock}{Objetivo Didáctico}
		\footnotesize
		Formalizar el paso de espacios muestrales abstractos a funciones de probabilidad discretas ($\sum f(x)=1$), estableciendo la noción de soporte medible y caracterizando la masa asignada a cada evento del fenómeno aleatorio.
	\end{alertblock}
\end{frame}

% SLIDE 03: Motivación I
\begin{frame}{Motivación: Del Espacio Muestral Abstracto a los Números Reales}
	\begin{columns}[T]
		\begin{column}{0.48\textwidth}
			\begin{block}{Limitación del Espacio Muestral $\Omega$}
				\footnotesize
				En fenómenos aleatorios complejos, el espacio muestral elemental $\Omega$ puede estar constituido por categorías no numéricas:
				\begin{itemize}
					\item Secuencias de monedas ($\set{AAS, SAA, \dots}$).
					\item Defectos cualitativos en manufactura ($\set{\text{Óptimo}, \text{Falla}}$).
					\item Perfiles de clientes o historiales clínicos.
				\end{itemize}
			\end{block}
		\end{column}
		\begin{column}{0.48\textwidth}
			\begin{block}{La Necesidad de Cuantificar}
				\footnotesize
				Para aplicar el poderoso herramental del análisis matemático y el álgebra lineal, es indispensable **proyectar** cada resultado cualitativo $\omega \in \Omega$ hacia un valor numérico unívoco en la recta real $\mathbb{R}$.
				
				\vspace{0.05cm}
				Esta proyección codifica cuantitativamente la propiedad empírica u operativa de interés para el ingeniero o científico de datos.
			\end{block}
		\end{column}
	\end{columns}
	\vspace{-0.2cm}
	\begin{alertblock}{Principio Fundamental de Modelación}
		\footnotesize
		Una \emph{variable aleatoria} es el puente formal entre el espacio de probabilidad abstracto $(\Omega, \mathcal{F}, P)$ y el análisis cuantitativo sobre $\mathbb{R}$.
	\end{alertblock}
\end{frame}

% SLIDE 04: Motivación II
\begin{frame}{Motivación: Proyección Algebraica sobre el Soporte}
	\begin{block}{Representación Algebraica de la Proyección}
		\small
		Si lanzamos tres monedas al aire, el espacio elemental posee $2^3 = 8$ puntos:
		\begin{align*}
			\Omega = \set{AAA, AAS, ASA, SAA, ASS, SAS, SSA, SSS}.
		\end{align*}
		Si definimos $X(\omega)$ como el número total de águilas ($A$), la función asigna:
		\begin{align*}
			X(AAA) = 3, \quad X(AAS) = X(ASA) = X(SAA) = 2, \dots
		\end{align*}
	\end{block}
	\vspace{-0.15cm}
	\begin{columns}[T]
		\begin{column}{0.48\textwidth}
			\begin{block}{Reducción de Dimensionalidad}
				\footnotesize
				El estudio de los 8 resultados abstractos se reduce a analizar únicamente 4 enteros discretos: $\set{0, 1, 2, 3}$.
			\end{block}
		\end{column}
		\begin{column}{0.48\textwidth}
			\begin{block}{Concentración de Masa}
				\footnotesize
				La probabilidad del espacio muestral se concentra y distribuye matemáticamente sobre los puntos del soporte discreto.
			\end{block}
		\end{column}
	\end{columns}
\end{frame}

% SLIDE 05: Definición Formal de Variable Aleatoria
\begin{frame}{Definición Formal: Variable Aleatoria Discreta}
	\begin{block}{Definición de Variable Aleatoria}
		\small
		Una \textbf{variable aleatoria} (o \emph{función estocástica}) es una función medible $X: \Omega \to \mathbb{R}$ que asigna a cada punto elemental $\omega \in \Omega$ un número real único $X(\omega) = x$.
	\end{block}
	\vspace{-0.1cm}
	\begin{block}{Caracterización Discreta}
		\small
		Se dice que $X$ es una \textbf{variable aleatoria discreta} si el rango o conjunto de valores posibles que puede asumir es finito o infinito numerable. Dicho conjunto se denomina \textbf{soporte} de $X$, denotado como:
		\begin{align*}
			\mathcal{S}_X = \set{x_1, x_2, \dots, x_k, \dots} \subset \mathbb{R}.
		\end{align*}
	\end{block}
	\vspace{-0.15cm}
	\begin{alertblock}{Convención Notacional}
		\footnotesize
		Las variables aleatorias se denotan rigurosamente con letras mayúsculas ($X, Y, Z$), mientras que sus realizaciones particulares se denotan con minúsculas ($x, y, z$).
	\end{alertblock}
\end{frame}

% SLIDE 06: Soporte de una Variable Aleatoria
\begin{frame}{Soporte Medible y cardinalidad Discreta}
	\begin{columns}[T]
		\begin{column}{0.48\textwidth}
			\begin{block}{Soporte Finito}
				\small
				El conjunto $\mathcal{S}_X$ consta de un número acotado de elementos ($|\mathcal{S}_X| = n < \infty$).
				\begin{itemize}
					\item \textbf{Suma de dos dados:} $\mathcal{S}_X = \set{2, 3, \dots, 12}$ con cardinalidad $11$.
					\item \textbf{Defectos en un lote:} $\mathcal{S}_X = \set{0, 1, \dots, N}$.
				\end{itemize}
			\end{block}
		\end{column}
		\begin{column}{0.48\textwidth}
			\begin{block}{Soporte Infinito Numerable}
				\small
				El conjunto $\mathcal{S}_X$ es biyectivo con los enteros naturales ($\mathbb{N}$ o $\mathbb{Z}^+$).
				\begin{itemize}
					\item \textbf{Ensayos hasta primer éxito:} $\mathcal{S}_X = \set{1, 2, 3, \dots}$ (Geométrica).
					\item \textbf{Llegadas a un servidor:} $\mathcal{S}_X = \set{0, 1, 2, \dots}$ (Poisson).
				\end{itemize}
			\end{block}
		\end{column}
	\end{columns}
	\vspace{-0.15cm}
	\begin{block}{Exclusión de Puntos Nulos}
		\footnotesize
		Por definición estricta, el soporte medible de una variable discreta excluye todo número real que posea una probabilidad idéntica a cero: $\mathcal{S}_X = \set{x \in \mathbb{R} : P(X = x) > 0}$.
	\end{block}
\end{frame}

% SLIDE 07: Función de Masa de Probabilidad (PMF)
\begin{frame}{Función de Masa de Probabilidad (PMF)}
	\begin{block}{Definición Formal de la PMF}
		\footnotesize
		Sea $X$ una variable aleatoria discreta con soporte $\mathcal{S}_X$. La **Función de Masa de Probabilidad** (o *función de probabilidad discreta*), denotada por $p_X(x)$ o simplemente $f(x)$, es la regla que asigna a cada número real $x$ la probabilidad exacta de que $X$ tome dicho valor:
		\begin{align*}
			f(x) = P(X = x) = P\left(\set{\omega \in \Omega : X(\omega) = x}\right).
		\end{align*}
	\end{block}
	\vspace{-0.2cm}
	\begin{block}{Extensión a toda la Recta Real $\mathbb{R}$}
		\footnotesize
		De manera exhaustiva, la función se define sobre todo $\mathbb{R}$ mediante la pieza nula fuera del soporte:
		\begin{align*}
			f(x) = \begin{cases}
				P(X = x_k) & \text{si } x = x_k \in \mathcal{S}_X, \\[2pt]
				0 & \text{en cualquier otro caso ($x \notin \mathcal{S}_X$).}
			\end{cases}
		\end{align*}
	\end{block}
\end{frame}

% SLIDE 08: Axiomas de Kolmogorov y Normalización
\begin{frame}{Axiomas de Kolmogorov para Funciones Discretas}
	\begin{block}{Teorema de Existencia y Legitimidad de una PMF}
		\small
		Una función real $f(x)$ constituye una función de masa de probabilidad legítima si y sólo si satisface rigurosamente los siguientes dos axiomas de Kolmogorov:
		\begin{enumerate}
			\itemsep=0.15cm
			\item \textbf{No negatividad:} Para todo número real $x \in \mathbb{R}$, se cumple que
			\begin{align*}
				f(x) \ge 0.
			\end{align*}
			\item \textbf{Normalización (Suma unitaria):} La suma total de las probabilidades asignadas a los puntos de su soporte medible es exactamente igual a la unidad:
			\begin{align*}
				\sum_{x \in \mathcal{S}_X} f(x) = 1.
			\end{align*}
		\end{enumerate}
	\end{block}
\end{frame}

% SLIDE 09: Ejemplos Canónicos del Textbook I
\begin{frame}{Ejemplos Canónicos del Libro: Monedas y Dados}
	\begin{columns}[T]
		\begin{column}{0.48\textwidth}
			\begin{block}{Ejemplo 2.6.2: Dos Monedas ($X = \text{Soles}$)}
				\footnotesize
				Espacio $\Omega = \set{HH, HT, TH, TT}$ equiprobable.
				Soporte $\mathcal{S}_X = \set{0, 1, 2}$.
				\begin{align*}
					f(x) = \begin{cases}
						1/4 & x = 0 \text{ ($TT$)} \\
						1/2 & x = 1 \text{ ($HT, TH$)} \\
						1/4 & x = 2 \text{ ($HH$)} \\
						0 & \text{en otro caso.}
					\end{cases}
				\end{align*}
			\end{block}
		\end{column}
		\begin{column}{0.48\textwidth}
			\begin{block}{Ejemplo 2.6.3: Suma de Dos Dados}
				\footnotesize
				Espacio $|\Omega| = 36$. Soporte $\mathcal{S}_X = \set{2, \dots, 12}$.
				Masa proporcional al número de formas de sumar $x$:
				\begin{align*}
					f(x) = \frac{6 - |x - 7|}{36}, \quad x \in \set{2, \dots, 12}.
				\end{align*}
				Por ejemplo: $f(7) = 6/36$, $f(2) = f(12) = 1/36$.
			\end{block}
		\end{column}
	\end{columns}
	\vspace{-0.15cm}
	\begin{alertblock}{Verificación Axiomática}
		\scriptsize
		En ambos casos se observa no negatividad estricta y aditividad unitaria ($\sum f(x) = \frac{1+2+1}{4} = 1$ y $\sum f(x) = \frac{36}{36} = 1$).
	\end{alertblock}
\end{frame}

% SLIDE 10: Ejemplos Canónicos del Textbook II
\begin{frame}{Ejemplos Canónicos del Libro: Familias con 3 Hijos}
	\begin{block}{Ejemplo 2.6.4: Distribución Binomial Elemental ($N=3, p=1/2$)}
		\small
		Sea $X$ el número de niños en una familia con 3 hijos, asumiendo independencia y equiprobabilidad de género. El espacio tiene $|\Omega| = 2^3 = 8$ puntos.
		
		\vspace{0.1cm}
		La función de masa se deriva por combinatoria sobre el soporte $\mathcal{S}_X = \set{0, 1, 2, 3}$:
		\begin{align*}
			f(0) = P(MMM) = \frac{1}{8}, \qquad & f(1) = P(NMM, MNM, MMN) = \frac{3}{8}, \\[4pt]
			f(2) = P(NNM, NMN, MNN) = \frac{3}{8}, \qquad & f(3) = P(NNN) = \frac{1}{8}.
		\end{align*}
	\end{block}
	\vspace{-0.2cm}
	\begin{center}
	\small
	\begin{tabular}{c|cccc|c}
		$x$ (Niños) & $0$ & $1$ & $2$ & $3$ & Suma \\
		\hline
		$f(x)$ & $0.125$ & $0.375$ & $0.375$ & $0.125$ & $1.000$
	\end{tabular}
	\end{center}
\end{frame}

% SLIDE 11: Lab Computacional I
\begin{frame}[fragile]{Laboratorio en Python --- Parte 1: Normalización de PMF Polinomial}
	\begin{block}{Comprobación Numérica de Constantes y Condiciones}
		\footnotesize
		Verificación de la constante de normalización $c$ para $f(x) = c x^2$ en $\mathcal{S}_X = \set{1, 2, 3, 4}$:
	\end{block}
	\vspace{-0.1cm}
	\lstinputlisting[language=Python, firstline=9, lastline=20, basicstyle=\fontsize{6.3pt}{7.5pt}\ttfamily]{../../code/03_variables_aleatorias_discretas/03.01_pmf_and_support.py}
\end{frame}

% SLIDE 12: Lab Computacional II
\begin{frame}[fragile]{Laboratorio en Python --- Parte 2: Probabilidad Condicional en PMF}
	\begin{block}{Evaluación de Subconjuntos Medibles en el Soporte}
		\footnotesize
		Cálculo programático del evento condicional $P(X \text{ es par} \mid X \ge 2)$ sobre la PMF normalizada:
	\end{block}
	\vspace{-0.1cm}
	\lstinputlisting[language=Python, firstline=21, lastline=27, basicstyle=\fontsize{6.5pt}{7.8pt}\ttfamily]{../../code/03_variables_aleatorias_discretas/03.01_pmf_and_support.py}
\end{frame}

% SLIDE 13: Lab Computacional III
\begin{frame}[fragile]{Laboratorio en Python --- Parte 3: Simulación Monte Carlo (Dados)}
	\begin{block}{Convergencia Empírica de la PMF de la Suma de Dos Dados}
		\footnotesize
		Simulación de $100,000$ lanzamientos de dados y comprobación del error absoluto máximo frente a la teoría:
	\end{block}
	\vspace{-0.1cm}
	\lstinputlisting[language=Python, firstline=29, lastline=44, basicstyle=\fontsize{6.0pt}{7.2pt}\ttfamily]{../../code/03_variables_aleatorias_discretas/03.01_pmf_and_support.py}
\end{frame}

% SLIDE 14: Lab Computacional IV
\begin{frame}[fragile]{Laboratorio en Python --- Parte 4: Transformación No Lineal de V.A.}
	\begin{block}{Inducción de Soporte y Proyección de Frecuencias}
		\footnotesize
		Cálculo automático de la PMF inducida por la transformación absoluta $Y = |X - 1|$ para $X$ uniforme:
	\end{block}
	\vspace{-0.1cm}
	\lstinputlisting[language=Python, firstline=46, lastline=60, basicstyle=\fontsize{6.3pt}{7.5pt}\ttfamily]{../../code/03_variables_aleatorias_discretas/03.01_pmf_and_support.py}
\end{frame}

% SLIDE 15: Síntesis de Problemas --- Nivel Fundamental
\begin{frame}{Síntesis de Problemas --- Nivel Fundamental (Sencillos / Directos)}
	\begin{columns}[T]
		\begin{column}{0.48\textwidth}
			\begin{block}{Problema 3.1.1: Normalización Polinomial}
				\small
				Dada $f(x) = c x^2$ en $\mathcal{S}_X = \set{1, 2, 3, 4}$, la condición de suma unitaria exige:
				\begin{align*}
					c(1 + 4 + 9 + 16) = 30c = 1 \implies c = \frac{1}{30}.
				\end{align*}
				La probabilidad condicional resulta:
				\begin{align*}
					P(X \text{ par} \mid X \ge 2) = \frac{20/30}{29/30} = \frac{20}{29}.
				\end{align*}
			\end{block}
		\end{column}
		\begin{column}{0.48\textwidth}
			\begin{block}{Problema 3.1.2: Tres Monedas y Diferencias}
				\small
				Si $X = |n_A - n_S|$ al lanzar 3 monedas, el soporte es $\mathcal{S}_X = \set{1, 3}$.
				\begin{align*}
					f(1) = \frac{6}{8} = \frac{3}{4}, \quad f(3) = \frac{2}{8} = \frac{1}{4}.
				\end{align*}
				Se comprueba fehacientemente que $3/4 + 1/4 = 1$.
			\end{block}
		\end{column}
	\end{columns}
\end{frame}

% SLIDE 16: Síntesis de Problemas --- Nivel Operativo
\begin{frame}{Síntesis de Problemas --- Nivel Operativo (Intermedios / Aplicados)}
	\begin{columns}[T]
		\begin{column}{0.48\textwidth}
			\begin{block}{Problema 3.1.5: Transformación Absoluta}
				\small
				Sea $X$ distribuida sobre $\set{-2, -1, 0, 1, 2}$ y $Y = |X - 1|$. El soporte inducido es $\mathcal{S}_Y = \set{0, 1, 2, 3}$.
				
				\vspace{0.05cm}
				Las masas se agrupan por preimágenes:
				\begin{align*}
					p_Y(1) = P(X=0) + P(X=2) = 0.30.
				\end{align*}
			\end{block}
		\end{column}
		\begin{column}{0.48\textwidth}
			\begin{block}{Problema 3.1.6: Máximo de Dos Esferas}
				\footnotesize
				Al extraer dos esferas sin reemplazo del $1$ al $6$, el máximo $M = \max(X_1, X_2)$ tiene masa:
				\begin{align*}
					f(m) = \frac{m-1}{15}, \quad m \in \set{2, \dots, 6}.
				\end{align*}
				La cola superior es $P(M \ge 5) = \frac{4+5}{15} = \frac{3}{5}$.
			\end{block}
		\end{column}
	\end{columns}
\end{frame}

% SLIDE 17: Síntesis de Problemas --- Nivel Analítico
\begin{frame}{Síntesis de Problemas --- Nivel Analítico (Avanzados / Demostraciones)}
	\begin{block}{Problema 3.1.7: Suma Telescópica y Normalización en $\mathcal{S}_N = \set{1, \dots, N}$}
		\footnotesize
		Dada la familia $f_N(x) = \dfrac{c_N}{x(x+1)}$, descomponemos mediante fracciones parciales:
		\begin{align*}
			\sum_{x=1}^{N} \frac{c_N}{x(x+1)} &= c_N \sum_{x=1}^{N} \left( \frac{1}{x} - \frac{1}{x+1} \right) \\
			&= c_N \left( 1 - \frac{1}{N+1} \right) = c_N \left( \frac{N}{N+1} \right) = 1.
		\end{align*}
		Por lo cual se demuestra rigurosamente que $c_N = \dfrac{N+1}{N}$.
	\end{block}
	\vspace{-0.15cm}
	\begin{block}{Problema 3.1.8: Transformación Cuadrática en Soporte Simétrico}
		\footnotesize
		Para $X$ uniforme sobre $\set{-m, \dots, m}$ ($2m+1$ puntos), la variable $Z = X^2$ concentra masa doble en los enteros positivos ($p_Z(k^2) = \frac{2}{2m+1}$ por adición de $\pm k$) y masa simple en el origen ($p_Z(0) = \frac{1}{2m+1}$).
	\end{block}
\end{frame}

% SLIDE 18: Síntesis de Problemas --- Nivel Desafiante
\begin{frame}{Síntesis de Problemas --- Nivel Desafiante (Expertos / Paradojas)}
	\begin{columns}[T]
		\begin{column}{0.48\textwidth}
			\begin{block}{Problema 3.1.9: Ley de Zipf y Colas de Pareto}
				\small
				Para $f(k) = \dfrac{1}{\zeta(s) k^s}$ en $\mathbb{Z}^+$, la suma unitaria se garantiza por la convergencia de la función zeta $\zeta(s)$.
				
				\vspace{0.05cm}
				La acotación integral demuestra un decaimiento polinomial (larga cola):
				\begin{align*}
					P(X > M) \sim \frac{1}{(s-1) M^{s-1}}.
				\end{align*}
			\end{block}
		\end{column}
		\begin{column}{0.48\textwidth}
			\begin{block}{Problema 3.1.10: Redundancia Discreta}
				\small
				En dos subsistemas independientes con probabilidades de falla $p_A, p_B$, el tiempo hasta la primera falla $T = \min(N_A, N_B)$ satisface:
				\begin{align*}
					P(T > n) = \left[ (1-p_A)(1-p_B) \right]^n.
				\end{align*}
				$T$ sigue una distribución geométrica exacta con parámetro $p_{\text{eq}} = p_A + p_B - p_A p_B$.
			\end{block}
		\end{column}
	\end{columns}
\end{frame}

% SLIDE 19: Conclusiones
\begin{frame}{Conclusiones y Resumen Conceptual}
	\begin{block}{Puntos Clave de la Sección 03.01}
		\small
		\begin{itemize}
			\itemsep=0.1cm
			\item \textbf{Variable Aleatoria Discreta:} Función medible que proyecta un espacio probabilístico sobre un soporte numérico finito o infinito numerable $\mathcal{S}_X$.
			\item \textbf{PMF Legítima:} Toda función de masa $f(x) = P(X = x)$ queda caracterizada de manera única por los axiomas $f(x) \ge 0$ y $\sum_{x \in \mathcal{S}_X} f(x) = 1$.
			\item \textbf{Transformaciones de Soporte:} Al aplicar una función $g(X)$, las probabilidades de las preimágenes disjuntas se suman algebraicamente para conformar la PMF inducida.
		\end{itemize}
	\end{block}
\end{frame}

% SLIDE 20: Puente Didáctico
\begin{frame}{Puente Didáctico hacia la Sección 03.02}
	\begin{block}{Próximo Paso: Función de Distribución Acumulada (CDF)}
		\small
		Una vez dominada la asignación puntual de probabilidad en cada entero del soporte mediante la PMF $f(x)$, estamos en posición de analizar las **probabilidades acumuladas** hasta un umbral dado:
		\begin{align*}
			F(x) = P(X \le x) = \sum_{u \le x} f(u).
		\end{align*}
		En la **Sección 03.02**, estudiaremos por qué la CDF de una variable discreta adopta invariablemente la geometría de una *función escalonada* continua por la derecha, y cómo reconstruir la PMF a partir de los saltos de dicha función.
	\end{block}
\end{frame}

\end{document}
